%% =====================================================================================
\section{Further results}\label{sec:further}
%% =====================================================================================
The comparison theorem applies wherever the distortion method does. We describe three applications: squarefree
moduli, odd moduli, and other rings. Full details are in separate notes accompanying this paper (the folders
\path{squarefree/}, \path{odd-moduli/} and \path{applications/}).

\subsection{Squarefree moduli}\label{sec:squarefree}

\emph{Status: new; verified.} The result below is certified by an outward-rounded computation, confirmed by an
independent exact recomputation, and formally verified in Lean~4 with Mathlib.

\begin{theorem}\label{thm:squarefree}
No finite set of congruences with distinct squarefree moduli, all at least $10$, covers $\Z$. Consequently every
covering system with distinct squarefree moduli has least modulus at most $7$.
\end{theorem}

Since $8$ and $9$ are not squarefree, moduli greater than $7$ are at least $10$, which gives the second statement.
The previous bound was $118$, due to Cummings, Filaseta and Trifonov \cite{CFT2025}. In the other direction,
Erd\H{o}s's covering
\[
  0(2),\ 0(3),\ 0(5),\ 1(6),\ 0(7),\ 1(10),\ 1(14),\ 2(15),\ 2(21),\ 23(30),\ 4(35),\ 5(42),\ 59(70),\ 104(105)
\]
has distinct squarefree moduli and least modulus $2$ (one checks directly that it covers $\Z/210$). Whether a
covering with distinct squarefree moduli and least modulus $3$ exists is a question of Selfridge that remains open
(see \cite[F13]{GuyUPINT} and \cite[Problem~4]{DaltonTrifonov2022}). Hence the largest least modulus of such a
covering lies in $\{2,3,5,6,7\}$, and by \cref{thm:sf7type} below in $\{2,3,5,6\}$.
% CHECK: edition and section of Guy's "Unsolved Problems in Number Theory" (cited as F13 in DELIV/squarefree/README.md).

\emph{Method.} For squarefree moduli every coordinate is $X_i=\Z/p_i$, every cylinder has depth $0$ or $1$ in each
coordinate, and the tilted variables enter only through the Bernoulli variables $\one\{V_i\ge1\}$, with parameter
$\rho_i=\min(1,\nu_i/p_i)$. The one-coordinate lemma then has a three-line proof.

\begin{lemma}\label{lem:bernoulli}
Let $K\ge0$ on $\Z/q$ with $\sum_yK(y)=1$ and $K\le\nu/q$, let $c\ge0$, let $w\ge0$ on $\Z/q$ with
$\sum_yw(y)=S$, and let $\varphi$ be convex and nondecreasing. Then, with $\rho=\min(1,\nu/q)$,
\[
  \sum_yK(y)\,\varphi\bigl(c+w(y)\bigr)\le(1-\rho)\,\varphi(c)+\rho\,\varphi(c+S).
\]
\end{lemma}

\begin{proof}
We may assume $S>0$. By convexity, $\varphi(c+w(y))\le(1-w(y)/S)\varphi(c)+(w(y)/S)\varphi(c+S)$. Moreover
$\sum_yK(y)w(y)/S\le\rho$, since $w(y)/S\le1$ and $K\le\nu/q$, and $\varphi(c+S)-\varphi(c)\ge0$.
\end{proof}

At level $p$ the moduli are $d=m'p$ with $m'$ squarefree and $(p-1)$-smooth, so the weights are $1/p$ rather than
the geometric weights of \cref{def:smooth}. The union bound, \cref{lem:pointwise}, the comparison theorem and the
enlargement give, for $0\le t\le\delta_p$,
\[
  P_p(B_p)\le\frac{\E\bigl[(N(s)/p-t)^+\bigr]}{1-t},\qquad N(s)=\#\{m'\mid s:\ m'p\ge m\},
\]
where $s=\prod_{q<p}q^{\epsilon_q}$ with independent $\epsilon_q\sim\mathrm{Bernoulli}(\rho_q)$. Let $L=\lceil m/p\rceil$
and let $s_A$ be the part of $s$ supported on the primes below $\min(p,L)$. A divisor with a prime factor at least $L$
is itself at least $L$, so
\[
  N(s)=2^{\omega(s)}-\#\{d\mid s_A:\ d<L\}
\]
exactly. The factor $s_A$ takes at most $2^{\pi(\sqrt m)}$ values and is enumerated, and $\omega(s/s_A)$ has a
Poisson-binomial law computed by dynamic programming, so the level bound is evaluated exactly, without the
decomposition of \cref{lem:AB}. Since $t\mapsto\E[(N/p-t)^+]/(1-t)$ is monotone between consecutive atoms of the law
of $N/p$, the best $t\in[0,\delta_p]$ is $0$, an atom below $\delta_p$, or $\delta_p$; the schedule takes each
$\delta_p$ either at an atom $2^k/p$ or equal to $\frac12$. Levels at which no modulus $\ge m$ can occur ($p\prod_{q<p}q<m$; for $m=10$ these are $p=2,3$)
are empty and are given $\nu_p=1$.

\emph{Certificate.} With $m=10$ and $X=\Pmax$, the certified total is $\eta_{\mathrm{up}}=\etaSqf$. The primes $5$
and $7$ contribute $0.208333$, the primes $11\le p<10^3$ contribute $0.745813$, those in $[10^3,10^4)$ contribute
$0.036571$, and the primes above $10^4$ less than $0.0067$. The tail passes the criterion of \cite[Theorem~6.1]{BBMST}:
$k=\kPrimes$, $\mu_k\ge0.0026745$ and $f_k\le5.1411\cdot10^8\le2.8386\cdot10^9$.
% RESOLVED 2026-09-27 (source files corrected): SQUAREFREE_PROOF.md states f_k <= 5.141e8, but the certificate line (certs/cert_m10_2e8.txt) has
%        f_gen_up = 514103195.66 > 5.141e8; the text uses 5.1411e8.
An independent program recomputed every level bound (as an exact rational for $p\le400$, as a $50$-digit interval for $p\le10^6$, and with its own
bound above) and found no violation; the exact sum of the level bounds up to $10^6$ is $0.997269049318$. Exact
simulations of the measures of BBMST on $509$ concrete squarefree systems found no violation of any step of the
chain, and the level bound is attained at some levels. As a further check, on the schedule of \cite{CFT2025} at
$m=119$ the same program reproduces the moment-bound totals of \cite{CFT2025}, while the hinge bound gives
$0.4567$ instead of $0.9990$. The same method proves the statement for moduli at least $11$ with a margin of about
$7\%$ ($\eta_{\mathrm{up}}=0.929935461718$).

\emph{Formal proof.} The Lean theorem \code{MinModulus.Squarefree.not\_covers} states \cref{thm:squarefree}
(moduli \code{d i} injective, squarefree and at least $10$), and it is complete. Besides \code{propext},
\code{Classical.choice} and \code{Quot.sound}, its only axiom is the auxiliary axiom introduced by
\code{native\_decide} for a numerical checker written in core Lean, which works in fixed point with scale $2^{128}$,
rounds upward, and certifies a total of $\etaSqfLean<1$. The link from the checker's output to the mathematical
statement is proved in Lean, and the mathematics reuses the formalized comparison theorem, distortion and tail
modules of \cref{sec:lean}.

\emph{Beyond $7$.} Least modulus at most $6$ (that is, $m=7$) is not proved: the per-level method gives
$\eta\approx1.091$, and the refinements we tried reach about $1.086$ (numerical). Splitting $\Z$ by the residue
modulo $30$ does much better numerically: if every class that the moduli dividing $30$ leave uncovered uses its own
schedule, chosen after the system is known, the worst systems we found by adversarial search still leave a class
with bound about $0.60$ for $m=7$, $0.66$--$0.68$ for $m=6$ and $0.81$--$0.84$ for $m=5$. We cannot yet certify
this. Weights and schedules fixed in advance cannot beat the joint bound, since an adversary that aligns every level
at an independent, uniformly random class forces the joint value; so a proof must branch on the system itself, as in
\cref{thm:sf23} below. The following partial results are rigorous (certified with outward rounding and recomputed in
exact rational arithmetic). For $(m,P_0)\in\{(3,19),(5,59),(6,113),(7,569)\}$, no system of congruences with
distinct, squarefree, $P_0$-smooth moduli, all at least $m$, covers $\Z$; here no tail argument is needed, and the
certified totals are $0.991683501684$, $0.988892548054$, $0.994080928716$ and $0.999184263699$. Combined with a
prime-replacement lemma (Simpson and Zeilberger \cite{SimpsonZeilberger1991}; Dalton and Trifonov
\cite[Lemma~13]{DaltonTrifonov2022}), which allows the largest prime $q\ge m$ used by a distinct squarefree covering
with moduli $\ge m$ to be replaced by any unused prime in $[m,q)$, this shows that such a covering uses at least
$102$, $28$ and $16$ primes $p\ge m$ for $m=7,6,5$ (for $m=7$ there is in fact no such covering, by
\cref{thm:sf7type} below, and for $m=6$ \cref{thm:sf6smooth} below raises $28$ to $45$). For $m=3$ the per-level method stops at $P_0=19$ (it gives about
$1.057$ at $P_0=23$), but a finer argument goes further.

\begin{theorem}\label{thm:sf23}
No finite set of congruences with distinct squarefree moduli, all at least $3$ and all $23$-smooth, covers $\Z$.
Consequently a covering system with distinct squarefree moduli and least modulus $3$ uses at least $9$ odd primes,
and at least $10$ primes in all; in particular the least common multiple of its moduli is at least
$2\cdot3\cdot5\cdots29=6469693230$.
\end{theorem}

\emph{Status: new; certified with outward rounding, rechecked by an independent program in exact rational
arithmetic (refereed by an AI agent), and formally verified in Lean~4 (\cref{sec:lean}).} We may assume every admissible modulus is used. The moduli $3,5,6,10,15,30$
leave an uncovered set $U\subseteq\Z/30$, and an enumeration of all $405000$ choices of their residues shows that $U$
contains the image, under the coordinatewise symmetries of $\Z/2\times\Z/3\times\Z/5$, of one of $7$
inclusion-minimal sets. For $z\in U$, a congruence modulo $ef$ with $e\mid30$ and $f>1$ meets the class
$z+30\Z$ exactly when $z$ lies in its class modulo $e$, and then acts on it as a progression modulo $f$; so the
class $z$ must be covered by a family of progressions in which each $f$ occurs at most $8$ times. The per-level
bound holds verbatim for families with repetitions (the comparison theorem does not use distinctness), and it shows
that the class is not covered when its bound is below $1$. A certificate consists of weights on pairs (class,
schedule) such that the weighted sum of the class bounds is below the total weight for every admissible pattern of
multiplicities; at each level and each $s$ the maximum over patterns is bounded by convexity, and a branch-and-bound
over the multiplicity patterns of the moduli $7e$ and $11e$, up to symmetry, closes all $7$ cases with $1905$
certificate records. The largest certified value is $0.999927670599$; a symmetry exchanging the coordinates modulo
$2$ and $3$ maps this case onto another one whose certificate gives $0.998954874683$. The corollary follows from the
prime-replacement lemma applied to odd primes and from \cite[Theorem~1.1]{BBMST2021ES}, which forces an even
modulus. The certificates, both checkers and a proof note are in \path{squarefree/beyond/smooth23/}.
Extending the branch-and-bound to the prime $29$ gives more.

\begin{theorem}\label{thm:sf29}
No finite set of congruences with distinct squarefree moduli, all at least $3$ and all $29$-smooth, covers $\Z$.
Consequently a covering system with distinct squarefree moduli and least modulus $3$ uses at least $10$ odd primes,
and at least $11$ primes in all; in particular the least common multiple of its moduli is at least
$2\cdot3\cdot5\cdots31=200560490130$.
\end{theorem}

\emph{Status: new; certified with outward rounding and recomputed by an independent program in exact rational
arithmetic (refereed by an AI agent); not formalized.} The proof is that of \cref{thm:sf23} with the levels up to
$29$. Five of the seven cases close with certificates of $311$ to $4050$ records. The other two, which are related by
a symmetry, need a much larger tree, which branches at every level up to $29$ and uses a richer menu of schedules,
with $187895$ records; its bounds were checked in parts, and the completeness of its tree on the whole certificate.
The largest certified value is $0.999999933977$. The corollary follows as before. The certificates, the checker and
a proof note are in \path{squarefree/beyond/smooth29/}.

For $m=7$ the same method, combined with a case split on the moduli $10$, $15$ and $30$, settles the question
completely, without any smoothness assumption.

\begin{theorem}\label{thm:sf7type}
No finite set of congruences with distinct squarefree moduli, all at least $7$, covers $\Z$. Consequently every
covering system with distinct squarefree moduli has least modulus at most $6$.
\end{theorem}

\emph{Status: new; certified with outward rounding and recomputed in full by an independent program in exact
rational arithmetic (refereed by an AI agent); formally verified in Lean~4 (\cref{sec:lean}).} The second statement
follows from the first, since
a covering with least modulus at least $7$ has all its moduli at least $7$. For a system with moduli at least $7$ let
$C\subseteq\Z/30$ be the set of classes covered by its congruences with moduli $10$, $15$ and $30$; it is contained
in a set $A\cup B\cup\{c\}$ with $A$ a class modulo $10$, $B$ a class modulo $15$ and $c$ a class modulo $30$. Adding
congruences preserves covering, so we may assume that all three moduli occur. Up to the coordinatewise symmetries of
$\Z/2\times\Z/3\times\Z/5$ the set $C$ is then contained in one of $9$ sets $C_1,\dots,C_9$ (an exhaustive check of
all $4500$ residue choices, and of the $956$ choices with some of the three moduli absent). The orbit of $C_9$ ($30$
sets of size $5$: a class modulo $5$ minus one point) covers the case $A\equiv B\pmod 5$; the orbits of
$C_1,\dots,C_8$ consist of sets of size $6$ with $A\not\equiv B\pmod 5$ and are distinguished by the class of $c$:
$C_6,C_7,C_8$ ($240$, $240$ and $720$ sets) have $c\not\equiv A\pmod 2$, $c\not\equiv B\pmod 3$ and $c\equiv A$,
$c\equiv B$ or neither modulo $5$; $C_1,C_3,C_5$ ($360$, $720$ and $360$ sets) have $c$ congruent to neither $A$ nor
$B$ modulo $5$, agreeing with $A$ modulo $2$ and with $B$ modulo $3$, with $A$ only, or with $B$ only; $C_4$ ($120$
sets) has $c\equiv A\pmod 5$ and $c\equiv B\pmod 3$; and $C_2$ ($240$ sets) has $c\equiv B\pmod 5$ and $c\equiv
A\pmod 2$. For each type the complement $U$ of $C_i$ in $\Z/30$ is alive, and each class $z\in U$ must be covered by
progressions modulo $f$ with multiplicities at most $8$, as before. The proof branches on the residues modulo $e$ of
the moduli $7e$ for $e\in\{2,3,5,6,10,15\}$ (up to the symmetries of $U$), maximizes exactly over the residue of
$210$, and bounds the levels $p\ge11$ with a common schedule $\delta_p=t_p=2^{k_p}/p$. For these levels the vertex
step of \cref{thm:sf23} is sharpened by applying Jensen's inequality across the groups of moduli with the same part
below $11$ and letting one profile serve all $s$ with the same part above $7$; with thresholds that are powers of
$2$ the resulting bound is a finite combination of exact profile maxima with universal coefficients, computed with
upward rounding. For type~$9$ ($25$ alive classes, $1480$ cells) this closes every cell, with largest value
$0.990728458807$; the coefficients are computed up to $2^{27}$ and the primes beyond are handled with
$\delta_p=\frac12$ and \cref{prop:tail}. For types $6$, $7$ and $8$ ($24$ alive classes) the cells whose bound is
not below $1$ ($1079$, $1356$ and $2110$ of $7272$, $7992$ and $16416$) are refined further by branching on the
residues modulo $e$ of the moduli $11e$, over all $27000$ choices: the free residue of $330$ and the modulus $77$
are treated either exactly or by a Jensen split, the levels $p\ge13$ by the same grouping over the parts in
$\{1,7,11,77\}$, and the primes beyond $2^{31}$ by the Cauchy--Schwarz inequality and \cref{prop:tail}. This
gives $15788$, $25842$ and $24929$ refined cells, all with certified value below $1$; the largest values are
$0.999791591902$, $0.999754401235$ and $0.999797679840$ before refinement and $0.999978915577$, $0.999979279565$
and $0.999979476961$ after. The other five types ($24$ alive classes) need one more level, and their schedule
differs from the common one at nine primes ($k_p$ is lowered by one at $37,71,137,257,263,491,499,947,953$; any
schedule with $2^{k_p}<p$ is admissible). Some cells refined at level $11$ are refined once more by branching on
the residues modulo $e$ of the moduli $13e$; there the moduli $91e$, $143e$ and $1001e$ are treated by a Jensen
split with two or four equal weights, the levels $p\ge17$ by the grouping over the eight parts dividing
$7\cdot11\cdot13$, and the primes beyond $2^{31}$ as before. At the hardest level-$13$ cells the Jensen split is
replaced by the exact expectation over $s\subseteq\{7,11\}$ at $p=13$, whose four terms are maximized separately
(the last one by the vertex step), the level-$13$ analogue of the exact treatment of $77$ at level $11$. In one cell
of type~$2$, $8$ level-$13$ cells remain above $1$ even so; each is split further according to the class modulo
$30$ of the congruence with modulus $390$, which may be taken to lie in $U$ (if the modulus is absent or its class is
dead, every one of the resulting bounds applies, since all terms are nondecreasing), and this gives $192$ cells, each
with its own weights. The numbers are as follows.
\begin{center}
\begin{tabular}{@{}rrrrrrl@{}}
\toprule
type & cells & refined & level-$11$ cells & level-$13$ cells & exact & largest value \\
\midrule
$3$ & $16416$ & $5396$ & $49365$ & $552$ & $0$ & $0.999979939775$ \\
$5$ & $8512$ & $2802$ & $60649$ & $1949$ & $0$ & $0.999979897908$ \\
$1$ & $8512$ & $4526$ & $574608$ & $61939$ & $307$ & $0.999979960$ \\
$4$ & $3744$ & $1345$ & $76189$ & $34010$ & $22$ & $0.999994881$ \\
$2$ & $7992$ & $3164$ & $270963$ & $92797$ & $2629$ & $0.999979968$ \\
\bottomrule
\end{tabular}
\end{center}
Here \emph{refined} is the number of cells refined at level $11$, the next two columns are the numbers of cells
produced at levels $11$ and $13$, and \emph{exact} is the number of level-$13$ cells that use the exact expectation
(for type~$2$ including the $192$ split cells); the largest values before refinement are between $0.99975$ and
$0.99978$. Every certificate was checked by the verifier with
upward rounding and recomputed in full by the referee in exact rational arithmetic; the certificates of types $1$ and
$2$ ($215$ and $127$ MB) were checked in pieces whose level-$7$ coverages are disjoint and together cover all
$27000$ choices. The certificates, verifiers and proof notes are in \path{squarefree/beyond/m7-type9/},
\path{squarefree/beyond/m7-types678/}, \path{squarefree/beyond/m7-type3/}, \path{squarefree/beyond/m7-type5/},
\path{squarefree/beyond/m7-type1/}, \path{squarefree/beyond/m7-type4/} and \path{squarefree/beyond/m7-type2/}.

Before the last types were settled, truncating the tail gave the weaker statement that no covering with distinct
squarefree moduli, all at least $7$ and all $1200$-smooth, exists, so that such a covering would involve at least
$194$ primes $p\ge7$ (with largest certified value $0.999638381232$); this is now subsumed by
\cref{thm:sf7type}, and its certificates remain in \path{squarefree/beyond/m7-smooth1200/}.

For moduli at least $6$ the block moduli are $6$, $10$, $15$ and $30$, and the same method gives a smooth analogue
of \cref{thm:sf7type}.

\begin{theorem}\label{thm:sf6smooth}
No finite set of congruences with distinct squarefree moduli, all at least $6$ and all $211$-smooth, covers $\Z$.
Consequently a covering system with distinct squarefree moduli, all at least $6$, uses at least $45$ primes $p\ge7$.
\end{theorem}

\emph{Status: new; certified with outward rounding and recomputed in full by an independent program in exact
rational arithmetic (refereed by an AI agent); not formalized.} The moduli $6$, $10$, $15$ and $30$ leave an uncovered
set $U\subseteq\Z/30$, and an exhaustive check of all $27000$ residue choices shows that, up to the coordinatewise
symmetries of $\Z/2\times\Z/3\times\Z/5$, $U$ contains one of $26$ inclusion-minimal sets, of sizes $19$, $20$ and
$21$ (so does every choice in which some of the four moduli are absent). Each type is closed by the cells of the
proof of \cref{thm:sf7type}, where now no level above $211$ occurs: the tail sums stop at $211$ and there is no
remainder term. Level-$7$ cells close $249834$ of the $251652$ cells, with largest value $0.999589971059$; the other
$1818$ are refined at level $11$ ($6291$ cells, $45$ of them with the exact treatment of $77$), and $16$ of these once
more at level $13$, with largest value $0.999978742660$. The corollary follows from the prime-replacement lemma, since
there are $44$ primes in $[7,211]$. Level-$7$ cells alone reach $167$ (so at least $37$ primes), and without a
smoothness assumption the hardest cell of every type stays above $1.1$ numerically even after all the refinements used
for \cref{thm:sf7type}; whether every covering with distinct squarefree moduli has least modulus at most $5$ remains
open. The certificates, the verifier with its smooth mode and the referee's report are in
\path{squarefree/beyond/m6-smooth211/}, and those for $167$ with level-$7$ cells alone in
\path{squarefree/beyond/m6-smooth167/}.

\subsection{Odd moduli and divisibility conditions}\label{sec:odd}

\emph{Status: new; certified with outward rounding and refereed by AI agents (two referee passes).
\Cref{thm:odd,thm:odd2} are also formally verified in Lean~4.}

In this subsection all moduli are greater than $1$. Erd\H{o}s and Selfridge asked whether a covering system with
distinct odd moduli exists; the question is open.

\begin{theorem}\label{thm:odd}
Let $\cA$ be a finite family of congruences with distinct moduli, all greater than $1$, that covers $\Z$.
\begin{enumerate}[label=(\Alph*)]
\item Some modulus is divisible by $2$ or by $9$.
\item Some modulus is divisible by $2$ or by $5$.
\item If all the moduli are odd, then some modulus is less than $11$.
\end{enumerate}
\end{theorem}

BBMST proved that $2\mid Q$, or $9\mid Q$, or $15\mid Q$, where $Q$ is the lcm of the moduli
\cite[Theorem~1.4]{BBMST}, and remark that in the last case they are unable to show that a single modulus is divisible
by $15$. Part~(A) shows that the third alternative is not needed. Part~(B) is the analogue, with $5$ in place of $3$,
of the theorem of Hough and Nielsen \cite{HoughNielsen2019} that some modulus is divisible by $2$ or $3$ (reproved in
\cite[Theorem~7.1]{BBMST}). By part~(C), a covering with distinct odd moduli, if one exists, has a modulus in
$\{3,5,7,9\}$.

\emph{Formal proof.} All three parts of \cref{thm:odd} are formally verified in Lean~4 with Mathlib, as the theorems
\code{MinModulus.Odd.not\_covers\_A}, \code{\_B} and \code{\_C} (\path{lean/MinModulus/Odd/}). Each depends only on
\code{propext}, \code{Classical.choice}, \code{Quot.sound} and one \code{native\_decide} evaluation of a numerical
checker written in Lean. The proof has three parts.
\begin{itemize}
\item A generic chain for moduli restricted to an admissible set. It uses the level bound with the enlargement to
  the admissible moduli, and the sharpened tail of \cref{prop:tailtwo}.
\item Per-case identities for the weights, which exclude $2$ (and $5$ in case~(B)) and cap $3$ at exponent $1$ in
  case~(A).
\item A checker that evaluates the resulting $\tau$-split bounds for the primes up to $2\cdot10^8$.
\end{itemize}
Its certified totals, tail included, are at most $0.998406$ for~(A), $0.784852$ for~(B) and $0.989342$ for~(C). Each native
evaluation takes about $80$ seconds.

\begin{theorem}\label{thm:odd2}
For each of the following six properties, every covering system with distinct moduli, all greater than $1$, has a
modulus with that property.
\begin{itemize}
\item[\textup{B(7)}] Divisible by $2$, or by $5r$ for some prime $r\ge7$.
\item[\textup{HN(13)}] Divisible by $2$, or by $3r$ for some prime $r\ge13$.
\item[\textup{T3(7)}] Even, equal to $3$, or divisible by $7$.
\item[\textup{T3(11)}] Even, equal to $3$, or divisible by $11$.
\item[\textup{HN$^*$(47)}] Even, equal to $3$, or divisible by $3r$ for some prime $r\ge47$.
\item[\textup{B$^*$(29)}] Even, equal to $3$, or divisible by $5r$ for some prime $r\ge29$.
\end{itemize}
\end{theorem}

In a covering with distinct odd moduli, HN(13) says that $3$ must occur together with a prime at least $13$ in some
modulus, which strengthens the theorem of Hough and Nielsen, and B(7) says that the moduli divisible by $5$ cannot all
be of the form $3^a5^b$. The statements HN($P$) for all primes $P$ together are equivalent to the conjecture that no
covering with distinct odd moduli exists, since such a covering violates HN($P$) for every $P$ larger than all its prime
factors; the same holds for the analogous statements B($P$). The last four statements of \cref{thm:odd2} show how much
of the difficulty comes from the modulus $3$ itself: once it is excluded, the method reaches much further.

\emph{Method.} In each case the hypothesis removes primes or prime powers from the construction of
\cref{sec:loss}. When no modulus is even, $2\nmid Q$, so the prime $2$ is neither a level nor a factor of $s$.
\begin{itemize}
\item (A) ($m=3$). Since $9\nmid Q$, the exponent of $3$ is capped at $1$: the variable $v_3$ takes the value $1$
  with probability $\nu_3/3$ and $0$ otherwise, level $3$ has $\alpha_3\le1/3$, and $U_p(s)=\tau(s)/(p-1)$ for the
  other levels. The tail uses \cite[Theorem~6.1]{BBMST} with $i_0=2$ and $\kappa=1+\nu_3$, as in the proof of
  \cite[Theorem~1.4]{BBMST}.
\item (B) ($m=3$). The prime $5$ is removed from the levels and from $s$; the tail uses $i_0=1$ and $\kappa=1$, as in
  all remaining cases.
\item (C) ($m=11$). Here $U_3=1/18$, $U_5(s)=(\tau(s)-\frac45)/4$, $U_7(s)=(\tau(s)-\frac67)/6$ and
  $U_p(s)=\tau(s)/(p-1)$ for $p\ge11$.
\item HN($P$) and B($P$). Let $q=3$ or $q=5$, respectively. The admissible moduli are the odd $d>1$ such that not both
  $q\mid d$ and the largest prime factor of $d$ is at least $P$. At the levels $p<P$ every odd modulus is admissible;
  at the levels $p\ge P$ no admissible modulus is divisible by $q$, so the prime $q$ is removed from $s$. The total is
  the sum of two certified runs, combined in exact rational arithmetic.
\item T3($p$), HN$^*$($P$) and B$^*$($P$). Among odd moduli greater than $1$, being different from $3$ is the same as
  being at least $5$, so these statements are proved with $m=5$ and the prime $2$ absent; in T3($p$) the prime $p$ is
  also removed, and HN$^*$($P$) and B$^*$($P$) are split into two runs as above.
\end{itemize}

\emph{Certificates.} The certified checker of \cref{sec:computation}, modified for these constraints, gives the totals
of \cref{tab:odd}; the tail criterion of \cite[Theorem~6.1]{BBMST}, with $X=\Pmax$ and $k=\kPrimes$, passes in each
case, against the threshold $(\log k+\log\log k-3)^2k\ge\critLo$. For (A)--(C), an independent referee reproduced the
certificates bit for bit and recomputed them with separately written code, obtaining totals of $0.998502$, $0.758524$
and $0.990483$ (for (A) with the fractional-moment bounds starting at $2\cdot10^6$). A second referee pass checked every
reduction of \cref{thm:odd2}, reproduced all $18$ runs and $5$ combined certificates bit for bit and re-certified every
statement with a newly written evaluator (for example $0.998865$ for HN$^*$(47) and $0.997346$ for B$^*$(29)). Neither
pass found errors. Exact simulations of the measures of BBMST on small systems with adversarial residues found no
violation of the certified level bounds.

\begin{table}[ht]
\centering
\caption{The certificates for \cref{thm:odd,thm:odd2}: the certified total $\eta_{\mathrm{up}}$ and an upper bound for
the quantity $f_k$ of the tail criterion, which must be at most $2.8386\cdot10^9$
(\code{odd-moduli/certs/}).}
\label{tab:odd}
\small
\begin{tabular}{@{}llr@{\qquad}llr@{}}
\toprule
statement & $\eta_{\mathrm{up}}$ & $f_k\le$ & statement & $\eta_{\mathrm{up}}$ & $f_k\le$\\
\midrule
(A) & $\etaOddA$ & $1.0648\cdot10^8$ & T3(7) & $0.730213115725$ & $7.389\cdot10^5$\\
(B) & $\etaOddB$ & $6.609\cdot10^5$ & T3(11) & $0.981833231256$ & $1.2194\cdot10^7$\\
(C) & $\etaOddC$ & $1.8175\cdot10^7$ & HN$^*$(47) & $0.998966442526$ & $2.8474\cdot10^8$\\
B(7) & $0.925376005735$ & $5.877\cdot10^6$ & B$^*$(29) & $0.997505909477$ & $1.1330\cdot10^8$\\
HN(13) & $0.947722737316$ & $6.609\cdot10^6$ & & &\\
\bottomrule
\end{tabular}
\end{table}
% Sources: DELIV/odd-moduli/README.md, ODD_MODULI_NOTE.md (sections 0, 4, 9, 10), certs/final_*.txt and
%          certs/fix_{A,B,C}.txt.  f_k: A 106471010.75, B 660860.43, C 18174206.68, B(7) 5.876307e6 (delta_5 = 1/2),
%          HN(13) 6.608906e6, T3(7) 474995 * 14/9 <= 738882, T3(11) 9237311 * 1.32 <= 1.2194e7, HN*(47) 2.847320e8,
%          B*(29) 1.132955e8.
% RESOLVED 2026-09-27 (source files corrected): ODD_MODULI_NOTE.md section 12 says that the six statements of Theorem 8.4 "have not had an independent referee
%        pass"; odd-moduli/README.md (later) reports a second referee pass with no errors.  The text follows the README.

\emph{Limits} (numerical; a total above $1$ means only that this method does not prove the statement). The method does
not prove that some modulus is divisible by $2$ or by $p$ for any prime $p\ge7$: the best totals found are $1.12$,
$1.55$ and $1.65$ for $p=7$, $11$ and $13$, and about $1.69$ for odd moduli without further restriction, that is, for
the Erd\H{o}s--Selfridge problem itself. Nor does it prove that a covering with distinct odd moduli has a modulus
below $9$ (best total $1.10$), so the threshold $11$ in (C) is the smallest that it certifies. The families stop
at HN(17) ($1.028$), HN$^*$(53) ($1.017$) and B$^*$(31) ($1.02$).
% RESOLVED 2026-09-27 (source files corrected): ODD_MODULI_NOTE.md section 10 lists 1.80 for p = 17 (B's schedule, fast setting), above the value 1.686 for
%        odd moduli without restriction, although the note and the README say that the values increase towards
%        about 1.69; the 1.80 is presumably an artifact of the unoptimized schedule.  The text omits p = 17.

\subsection{Function fields and imaginary quadratic fields}\label{sec:rings}

\emph{Status: new; certified with outward rounding and refereed by an AI agent with independently written checkers;
not formalized.} Details are in the notes \path{applications/fq_note.md} and \path{applications/nf_note.md}.

Let $R$ be a Dedekind domain in which every nonzero ideal $I$ has finite norm $N(I)=|R/I|$. A covering system of $R$
is a finite family of classes $a_i+I_i$, with nonzero proper ideals $I_i$, whose union is $R$; its moduli are distinct
if the ideals $I_i$ are pairwise distinct, and its multiplicity is the largest number of times an ideal occurs. For
$R=\mathbb{F}_q[x]$ the moduli are identified with monic polynomials of positive degree.

\begin{theorem}\label{thm:rings}
Let $q$ be a prime power.
\begin{enumerate}[label=(\alph*)]
\item If $q\ge11$, no covering system of $\mathbb{F}_q[x]$ has distinct moduli.
\item Every covering system of $\mathbb{F}_q[x]$ with distinct moduli has a modulus of degree $1$ if $q\in\{8,9\}$,
  and a modulus of degree at most $2$ if $q=7$.
\item Let $K$ be a global function field of genus $0$ with constant field $\mathbb{F}_q$, let $S$ be a nonempty finite
  set of places of $K$, and let $O_S$ be the ring of $S$-integers. If $q\ge13$, no covering system of $O_S$ has
  distinct moduli.
\item If $q\ge29$, no covering system of $\mathbb{F}_q[x]$ has multiplicity at most $2$.
\end{enumerate}
Moreover, every covering system of $\Z[i]$ with distinct moduli has a modulus of norm at most $1109$, and every
covering system of $\Z[\omega]$, $\omega=(1+\sqrt{-3})/2$, with distinct moduli has a modulus of norm at most $99$.
\end{theorem}

Part (a) was known for $q\ge759$ \cite{LWWY2025} and then for $q>73$ \cite[Theorem~1.3]{BWang2026}; we certify every
prime power $11\le q\le79$. Part (b) answers, for $q=7,8,9$, the analogue of the minimum modulus problem raised by Li,
Wang, Wang and Yi \cite{LWWY2024}; together with (a) it gives an explicit bound for the least degree of a modulus for
every $q\ge7$. (For $q=2$ a covering with distinct moduli exists, for example $0\bmod x$, $0\bmod x+1$,
$1\bmod x^2+x$; the cases $q\in\{3,4,5\}$ are open.) Part (c) follows from \cite[Theorem~1.1]{BWang2026}, with genus
$0$ and multiplicity $1$, for $q\ge83$; we certify every prime power $13\le q\le81$. For part (d),
\cite[Theorem~1.1]{BWang2026} with multiplicity $2$ requires $q>4\cdot82.26=329.04$, that is, $q\ge331$; we certify all
$68$ prime powers $29\le q\le317$, with two independent checkers, and there is no prime power between $317$ and $331$.
For the last statement, Li, Wang and Yi \cite{LiWangYi2026} proved that such a bound exists for every number field,
with a constant that is not explicit; as far as we know, these are the first explicit minimum-norm bounds for distinct
covering systems of a number field other than $\mathbb{Q}$.

\emph{Why the method transfers.} The arguments of \cref{sec:distortion,sec:comparison,sec:loss} use only the
following structure.
\begin{enumerate}[label=(\arabic*)]
\item \emph{Order and CRT.} Fix a total order on the nonzero prime ideals compatible with the norm (for function fields,
  with the degree), with ties broken arbitrarily but once and for all. If $Q=\bigcap_iI_i=\prod_j\mathfrak p_j^{\gamma_j}$,
  then $R/Q\cong\prod_jR/\mathfrak p_j^{\gamma_j}$, and every class $a+I$ with $I\mid Q$ is a product of balls
  $a+\mathfrak p_j^{v}$; a ball of depth $v$ in $R/\mathfrak p^{\gamma}$ has $N(\mathfrak p)^{\gamma-v}$ elements.
  The distortion method runs level by level along this order, and \cref{lem:F1,lem:F2,lem:F3,prop:F4,lem:kernel}
  hold with the same proofs.
\item \emph{Comparison.} \Cref{thm:comparison} holds verbatim with $X_i=R/\mathfrak p_i^{e_i}$ and $q_i=N(\mathfrak p_i)$:
  \cref{prop:rectangles} uses only the sizes of the sets, and \cref{lem:vlaw} only the sizes of the balls. (For the
  uniform measure, compare Kucheriaviy's analogue of Rogers' theorem in commutative rings \cite{Kucheriaviy2026}.)
\item \emph{The level bound.} At the level of $\mathfrak p$, a class of $B_{\mathfrak p}$ has modulus
  $\mathfrak g\mathfrak p^r$ with $r\ge1$ and $\mathfrak g$ composed of earlier primes, and the pair $(\mathfrak g,r)$
  determines the modulus by unique factorization. So \cref{lem:enlargement} holds, with a factor $\sigma$ for
  multiplicity $\sigma$, and the level bound is $\E[(U^{\mathrm{al}}-t)^+]/(1-t)$ with
  $U^{\mathrm{al}}=\sigma\sum_{\mathfrak g\mid\mathfrak s}w(\mathfrak g)$, where $\mathfrak s$ is a random ideal
  composed of the earlier primes with independent tilted exponents, $\P(v_{\mathfrak q}\ge a)=\min(1,\nu_{\mathfrak q}
  N(\mathfrak q)^{-a})$, and $w(\mathfrak g)=\sum N(\mathfrak p)^{-r}$ over the admissible $r\ge1$.
\item \emph{Ties.} Prime ideals of equal norm are distinct coordinates, distinct levels and distinct divisors; the
  order only decides at which level a class is counted, and the schedule depends only on the position in the order.
\item \emph{Dominance.} The level bound is nondecreasing in the set of earlier primes, and the schedule is fixed in
  advance, so the sum of the level bounds over all prime ideals of $R$ bounds the losses of every system. The
  checkers use the exact number of monic irreducible polynomials of each degree for $\mathbb{F}_q[x]$, the bound
  $\lfloor(q^n+1)/n\rfloor$ for the number of places of degree $n$ of a function field of genus $0$, and the lists of
  prime ideals of $\Z[i]$ and $\Z[\omega]$.
\item \emph{Tails.} For function fields, the degrees beyond the explicitly treated range are handled by
  second-moment bounds summed over all primes of each degree, with a geometric remainder. For $\Z[i]$ and
  $\Z[\omega]$, the prime ideals of norm above $X$ get $\delta=\frac12$, and an explicit estimate based on bounds of
  Rosser and Schoenfeld \cite{RosserSchoenfeld1962} for sums over primes replaces \cite[Theorem~6.1]{BBMST}.
\end{enumerate}

\emph{Certificates.} The checkers \path{applications/fqcert.c} and \path{applications/nfcert.c}, which follow the
rounding conventions of \cref{sec:rounding} (the second is derived from \code{rigcert.c}), certify the following
totals (rounded up): in (a), from $0.88877$ for $q=11$ down to $0.03433$ for $q=79$; in (b), $0.98782$ for $q=8$ and
$0.82650$ for $q=9$ (minimum degree $2$), and $0.73051$ for $q=7$ (minimum degree $3$); in (c), from $0.94213$ for $q=13$
down to $0.03918$ for $q=81$; in (d), from $0.99973$ for $q=29$ down to $0.02476$ for $q=317$; and for $\Z[i]$ and
$\Z[\omega]$, with minimum norms $1110$ and $100$, the sum $\eta_{\mathrm{up}}+\text{tail}$ is $0.994103491030+0.002750290372$ for $\Z[i]$
($X=2\cdot10^7$) and $0.998222202602+0.000610531031$ for $\Z[\omega]$ ($X=10^8$). The referee confirmed all five claims
with its own checkers: for (d) its $68$ totals agree with ours to within $2.6\cdot10^{-5}$, and for $\Z[i]$ and
$\Z[\omega]$ it obtained $0.993695$ and $0.997984$ and recomputed the tails in $60$-digit arithmetic. Brute-force and exact-rational checks at the
dominant levels agree as well. The method fails (numerically) for $q=7$ with minimum degree $2$ (total $1.21$) and for
$q=5$ with minimum degree $3$ ($1.45$); for $\Z[i]$ at minimum norm $1100$ and for $\Z[\omega]$ at $90$ the certified
totals are $1.0011$ and $1.0364$.
% Sources: DELIV/applications/PACKAGE.md, fq_note.md, nf_note.md and expected/claim_{A..E}.txt.
